%% kapa: exam=1 round=27 year=2016 period=2 form=A subject=appraisal_law subject_ko=감정평가관계법규 seq=36 number=36
%% source: Q-Net 감정평가사 기출문제 (cst00309) artlSeq=5169402 file=extracted/1008/5169402/제27회 감정평가사 1차 문제/제27회 감정평가사 2교시문제(A형).pdf
%% conversion: kordoc 4.24.51 -> tools/gen1.py (auto); review=proofread-sonnet; flags=none
\begin{question}{36}
동산ㆍ채권 등의 담보에 관한 법령상 동산담보권에 관한 설명으로 옳지 않은 것은?
\begin{choices}
  \item 창고증권이 작성된 동산은 담보등기의 목적물이 될 수 없다.
  \item 담보권설정자의 상호등기가 말소된 경우에도 이미 설정된 동산담보권의 효력에는 영향을 미치지 아니한다.
  \item 동산담보권은 피담보채권과 분리하여 타인에게 양도할 수 없다.
  \item 담보권자는 채권의 일부를 변제받은 경우에도 담보목적물 전부에 대하여 그 권리를 행사할 수 있다.
  \item 동산담보권의 효력은 법률에 다른 규정이 없거나 설정행위에 다른 약정이 없다면 담보목적물의 종물에 미치지 않는다.
\end{choices}
\end{question}
\answer{5}
